% Compile with XeLaTeX (the paper quotes Chinese throughout).
% Needs acl.sty and acl_natbib.bst from github.com/acl-org/acl-style-files
\documentclass[11pt]{article}
\usepackage[review]{acl}

\usepackage{times}
\usepackage{latexsym}
\usepackage{booktabs}
\usepackage{multirow}
\usepackage{graphicx}
\usepackage{microtype}
\usepackage[normalem]{ulem}

% Chinese. xeCJK needs XeLaTeX; change the font if it is unavailable on your
% system (Noto Serif CJK SC and Songti SC are common alternatives).
\usepackage{xeCJK}
% Songti SC ships with macOS and is a serif CJK face, which matches the Latin
% body text. On Overleaf or Linux, Noto Serif CJK SC is the equivalent and is
% usually present; swap the name if the build cannot find this one.
\setCJKmainfont{Songti SC}[AutoFakeBold=true]

% acl.sty loads hyperref itself, which is why \href and \url below need no
% package here. If the first build errors on an undefined \href, uncomment the
% next line -- but do NOT load it twice, that causes an option clash.
% \usepackage{hyperref}

\title{Which HSK 3.0? Two Official Documents, and Half of All\\
       Grading Decisions Change}

\author{Alvaro Serrano \\
  Independent researcher \\
  \texttt{alvaro.serrano.gp@gmail.com}}

\begin{document}
\maketitle

\begin{abstract}
``HSK 3.0'' names two different official documents.
《国际中文教育中文水平等级标准》 (GF0025-2021) is a national grading standard in
force since 1 July 2021; the 新版HSK考试大纲 is the examination syllabus
published in November 2025 and in force since July 2026. They are widely
conflated. We show the conflation is consequential. The two documents assign
different levels to \textbf{41.5\%} of the words and \textbf{40.7\%} of the
characters they share. Grading a corpus of 102 authentic graded readers against
one rather than the other changes the assigned level of \textbf{48\% of texts},
almost always upward, and shifts five of six shelf medians by a full level.
Measured against HSK 2.0, the two disagree about their own predecessor: 81.8\%
of jointly graded words moved under the 2021 standard against 51.1\% under the
2025 syllabus. We contribute a validated extraction of the 2025 syllabus,
previously available only as a 406-page PDF, whose per-level counts reproduce
the official totals exactly; an open, tested, dependency-free library that
grades against either document and reports which; an aligned 102-text
graded-reader corpus; and HSKBench, a controlled-difficulty generation
benchmark in which the grader is the metric rather than the source of labels.
Human authors writing to an explicit level target hit it 61.8\% of the time,
and two models 2.7 points apart in aggregate differ by 48 points on a single
level.
\end{abstract}

\section{Introduction}

Graded reading material is the backbone of second-language literacy
instruction, and grading it requires answering a narrow question: what does a
reader need to know to read this text? For Chinese, the answer has been
anchored for two decades to the HSK vocabulary lists.

That anchor moved, and it is not obvious \emph{to what}. Practitioners speak of
``HSK 3.0'' as a single thing. It is not. A national grading standard was
issued in 2021 \citep{moe2021standard} and an examination syllabus in 2025
\citep{clec2025syllabus}, four years apart, with different word lists,
different character inventories and different level boundaries. Popular
reference sites quote one document's level counts beside the other's
vocabulary; tools report ``HSK 3.0'' levels without saying which document
produced them.

We show this is not a pedantic distinction. Nearly half of real texts receive a
different level depending on which document is used, and the disagreement runs
in a consistent direction. A learner placed by a tool calibrated to the 2021
standard is being placed against a document that is not the examination they
will sit.

\noindent Contributions:

\begin{enumerate}
\item \textbf{A validated extraction of the 2025 examination syllabus}
  (\S\ref{sec:extract}). Our per-level entry counts reproduce the official
  cumulative totals exactly, and we find it grades \textbf{3{,}088} recognition
  characters, correcting the 3{,}079 widely reported by secondary sources.
\item \textbf{A quantification of how far the documents disagree}
  (\S\ref{sec:disagree}). The syllabus's own development paper
  \citep{feng2026hsk30} reports overlap but not level change; we distinguish
  the two, and show that high overlap coexists with high disagreement.
\item \textbf{A measurement of what that costs} (\S\ref{sec:cost}): 48\% of
  texts in an authentic graded-reader corpus change level.
\item \textbf{A character-level grading method and library}
  (\S\ref{sec:method}) that grades against either document and records which.
\item \textbf{A corpus and a benchmark} (\S\ref{sec:corpus},
  \S\ref{sec:bench}), the latter measuring difficulty-\emph{controlled
  generation}, on which human authors score 61.8\%.
\end{enumerate}

\section{Background}

\subsection{Three documents, routinely conflated}

\begin{table}[t]
\centering\small
\begin{tabular}{@{}llr@{}}
\toprule
\textbf{Document} & \textbf{Date} & \textbf{Words} \\
\midrule
HSK 2.0 exam lists & 2009--10 & 4{,}991 \\
GF0025-2021 & 1 Jul 2021 & 10{,}916 \\
2025 exam syllabus & Jul 2026 & 10{,}896 \\
\bottomrule
\end{tabular}
\caption{What ``HSK 3.0'' can mean.}
\label{tab:documents}
\end{table}

The 2021 document is a 语言文字规范 issued by the Ministry of Education and the
State Language Commission, covering syllables, characters, vocabulary and
grammar across three levels and nine bands. It is the document with public
machine-readable transcriptions, and therefore the one most existing tooling
uses --- including, until this work, our own.

The 2025 document is the examination syllabus, issued by the Center for
Language Education and Cooperation (中外语言交流合作中心), and governs the
examination itself.

Both grade characters separately from words, which matters: the character
inventory is what gates reading (\S\ref{sec:charlevel}). Levels 7, 8 and 9 are
a single undifferentiated band in both; tools advertising an ``HSK 8 word
list'' have invented a split neither document makes.

\subsection{Coverage thresholds}

\citet{hu2000unknown} found that around 98\% coverage of running words was
needed for adequate unassisted comprehension. \citet{laufer2010lexical}
proposed two thresholds: an optimal 98\% and a \textbf{minimal 95\%}, at which
comprehension is adequate \emph{with support}. Graded readers are read with
support, so we adopt 95\% as the default and expose it as a parameter. The 98\%
figure has since been partially challenged: \citet{kremmel2023replicating}
could not fully replicate it in a different population, and the original
regression rested on 66 participants.

This literature measures coverage over \emph{words}, in English. We apply the
threshold to Chinese characters for the reasons in \S\ref{sec:charlevel}; that
transfer is a principled assumption, not a result.

\section{Extracting the 2025 syllabus}
\label{sec:extract}

The syllabus is published as a 406-page PDF with a genuine text layer, so
extraction is parsing rather than OCR. Vocabulary occupies pp.~79--354 as rows
of 序号 等级 词语 拼音 词性; recognition characters (认读字) occupy
pp.~356--376, numbered per level.

Two details are easy to get wrong. \textbf{The level field carries more than a
level}: a row reading 1（4） means the word is introduced at level 1 and has a
further sense graded at level 4. We take the primary level. \textbf{The
published totals count entries, not words}: the syllabus numbers 11{,}000
entries, but homographs graded at two levels appear as separate rows
(所 / 所2). These collapse to \textbf{10{,}896 distinct words}.

Validation is exact: parsed entry counts per level reproduce the official
cumulative figures (300 / 500 / 1{,}000 / 2{,}000 / 3{,}600 / 5{,}400 /
11{,}000). For characters, each level's parsed count equals that level's
maximum item index, with no cross-level duplicates, giving \textbf{3{,}088}
recognition characters. Secondary sources widely report 3{,}079; we believe
that figure to be incorrect.

\section{How far the documents disagree}
\label{sec:disagree}

\begin{table}[t]
\centering\small
\begin{tabular}{@{}lrrr@{}}
\toprule
\textbf{Comparison} & \textbf{Shared} & \textbf{Same} & \textbf{Moved} \\
\midrule
2.0 $\rightarrow$ 2021    & 4{,}482 & 814     & 81.8\% \\
2.0 $\rightarrow$ 2025    & 4{,}802 & 2{,}349 & 51.1\% \\
2021 $\rightarrow$ 2025   & 9{,}674 & 5{,}662 & \textbf{41.5\%} \\
\bottomrule
\end{tabular}
\caption{Pairwise level agreement.}
\label{tab:agreement}
\end{table}

\paragraph{Membership is not agreement.} The syllabus's developers report a
99.2\% vocabulary overlap rate with HSK 2.0 and describe the revision as
optimisation rather than replacement \citep{feng2026hsk30}. Both
characterisations concern \emph{membership} --- whether a word appears at all.
Our membership figures are lower (96.2\% of our HSK 2.0 list appears in the
syllabus), plausibly because our HSK 2.0 list is a public reconstruction of
4{,}991 words rather than the official 5{,}000-word 词汇大纲, and because the
published figure's computation is not stated. We do not dispute it. The point
is that a high overlap rate and a high level-disagreement rate are compatible,
and only the second determines whether a text is graded correctly. The words
are largely the same words; they are not at the same levels.

\paragraph{The two ``HSK 3.0'' documents disagree with each other} on 41.5\% of
shared vocabulary. They are not drafts of one another; they are separate
gradings.

\paragraph{They also disagree about their predecessor.} Judged against the 2021
standard, HSK 2.0 looks almost entirely regraded (81.8\% moved). Judged against
the syllabus, barely half moved. A practitioner who read ``81.8\% of HSK
vocabulary changed level'' and acted on it would have substantially
overestimated the disruption to the actual examination.

\paragraph{Characters diverge similarly.} The 2021 standard grades exactly
3{,}000 characters (300 per level 1--6, 1{,}200 at 7--9); the syllabus grades
3{,}088 with a markedly uneven distribution (246 / 125 / 284 / 441 / 431 / 413
/ 1{,}148). Of the 2{,}945 characters both grade, \textbf{1{,}200 (40.7\%) sit
at different levels}.

Within the HSK 2.0 $\rightarrow$ 2021 comparison the instability is
concentrated above the beginner band: 91\% of HSK 2.0 level-1 words keep their
level, against 35\% at level 2 and 11--23\% from level 3 upward. The beginner
foundation is stable under every comparison, which may explain why the scale of
the change has attracted little attention.

\section{What it costs in practice}
\label{sec:cost}

We regraded all 102 texts of our corpus (\S\ref{sec:corpus}) under both
documents, holding method, threshold and proper-noun policy constant so the
document is the only variable. \textbf{49 of 102 texts (48.0\%) change level};
43 become harder and 6 easier. Shifts are almost all a single level.

\begin{table}[t]
\centering\small
\begin{tabular}{@{}lcc@{}}
\toprule
\textbf{Shelf} & \textbf{2021} & \textbf{2025} \\
\midrule
newbie       & HSK 2 & HSK 3 \\
beginner     & HSK 2 & HSK 3 \\
intermediate & HSK 3 & HSK 4 \\
upper        & HSK 4 & HSK 5 \\
advanced     & HSK 4 & HSK 5 \\
native       & HSK 5 & HSK 5 \\
\bottomrule
\end{tabular}
\caption{Median measured level per shelf. Five of six shift a full level.}
\label{tab:shelves}
\end{table}

The mechanism is simple. The syllabus grades \textbf{371 characters at levels
1--2 combined, against the 2021 standard's 600} --- a beginner inventory 38\%
smaller. A text reaching 95\% coverage within 600 characters frequently cannot
reach it within 371, so it lands a level higher.

This is the \emph{same fact} the developers describe as ``significantly
reducing the vocabulary difficulty of the entry stage''
(显著降低了入门阶段的词汇难度; \citealp{feng2026hsk30}), viewed from the other
side. Less is demanded of a beginning learner, which is a pedagogical
improvement; the arithmetic consequence is that beginner material grades
higher. Both statements are true, and a tool reporting levels needs to model
the second while its users read about the first.

\textbf{Any tool reporting an ``HSK 3.0 level'' should state which document it
used.} Ours records it on every result.

\subsection{Robustness}

Two obvious objections are that 48\% reflects a convenient choice of the 95\%
bar, and that it reflects something peculiar to these 102 texts. Neither holds.
Across thresholds from 0.80 to 1.00, between 43\% and 62\% of texts change
level, and the 48\% we report at 0.95 is among the lowest values in the sweep.

The \emph{direction} is threshold-dependent and reverses: below 0.95 the
syllabus grades texts harder almost without exception, while at 0.98 and above
the balance tips the other way. The syllabus allocates fewer characters to the
beginner levels but 88 more in total, so the smaller beginner inventory
dominates at low bars and the larger total inventory helps as the bar
approaches full coverage.

We also replicated on a \textbf{disjoint set of 30 texts} from a separate
content stream, which played no part in any decision reported here:
\textbf{14 of 30 changed level (46.7\%), all of them harder}. The finding is
not a property of the released corpus.

\section{Method}
\label{sec:method}

\subsection{Character-level, not word-level}
\label{sec:charlevel}

Word-level grading fails on segmented Chinese for a mundane reason: segmenters
emit tokens absent from any graded list. Phrase tokens such as 我的, 七点 and
蓝色 are natural segmentation output but not list entries, so the ``unknown''
share exceeds 95\% at every level and ordinary beginner text grades as
off-scale. Both documents' separate character gradings resolve this, and
reflect a fact about Chinese: the character inventory, not the word inventory,
gates reading.

\subsection{The official character list, not a derived one}

Deriving a character's level from the lowest-level word containing it agrees
with the 2021 official list on 2{,}962 of 2{,}969 derivable characters --- and
misgrades seven: 哥, 妈, 妹, 弟, 王, 第, 零. The failure is systematic: 哥, 妈,
妹 and 弟 are level-1 characters whose only listed words are the reduplicated
哥哥 and 妈妈, graded at level 4. Derivation misgrades precisely the vocabulary
a beginner meets first.

\subsection{Proper nouns}

Names should not count towards difficulty: a reader does not need a character's
name in productive vocabulary, because it is glossed in place. In our corpus a
text titled \emph{My Puppy Doudou} graded two levels high entirely on 豆豆, the
dog's name.

Proper nouns are romanised with a leading capital, making them detectable where
pinyin is available. One trap: an ASCII \verb|^[A-Z]| test silently fails on
accented capitals --- Ōuzhōu (欧洲), Ōuyà (欧亚), Ā Q Zhèngzhuàn (阿Q正传).
Since 欧 and 洲 are both advanced-band characters, missing one place name moved
a text two levels.

\subsection{Thresholds and character budgets}

A text's level is the lowest level at which cumulative character coverage
reaches the threshold. Characters outside the graded inventory count in the
denominator but never accumulate.

One consequence is not obvious to authors: reaching a 95\% bar means holding
the above-target share below 5\%, so \textbf{any single character exceeding
that budget blocks the target alone}. In a 75-character passage a fourth
repetition of one advanced character suffices --- we observed a stylistic edit
repeating 糕 a fourth time silently regress a text from level 3 to level 5.

\subsection{Grading collections}

Pooling a shelf's characters lets a few hard texts speak for all: pooling
reported ``HSK 3'' for the beginner shelf, on which 13 of 22 texts individually read at
HSK 1--2. We report the \textbf{median text}, and the interquartile range
rather than a wider window.

\subsection{Reading is not writing, and the gap is a constant}
\label{sec:writing}

HSK 3.0 grades characters twice: 认读字 (recognition, 3,088 under the 2025
syllabus) gates reading, while 书写字 (writing, 1,200) is what a learner must
produce by hand. The writing list is a strict subset of the recognition list.
No machine-readable version of it existed before this work.

Applying a coverage threshold to the writing list directly does not work, and
the failure is total: only 1,200 characters are graded for writing, so a median
text has around 60\% of its characters in the writing curriculum at all, and
\textbf{every text in our corpus misses a 95\% bar at every level}. We report
two numbers instead --- the \textbf{ceiling}, the share writable at any level,
and a \textbf{level} computed over the writable subset alone.

\begin{table}[t]
\centering\small
\begin{tabular}{@{}lccr@{}}
\toprule
\textbf{Shelf} & \textbf{Read} & \textbf{Write} & \textbf{Ceiling} \\
\midrule
newbie       & 2 & 3 & 58.2\% \\
beginner     & 3 & 3 & 57.9\% \\
intermediate & 4 & 4 & 60.5\% \\
upper        & 4 & 4 & 57.7\% \\
advanced     & 5 & 5 & 62.5\% \\
native       & 5 & 5 & 61.8\% \\
\bottomrule
\end{tabular}
\caption{Median reading level, writing level and writable ceiling per shelf.}
\label{tab:writing}
\end{table}

\paragraph{Writing difficulty tracks reading difficulty.} Median writing and
reading levels agree on five of six shelves, and the mean per-text gap stays
between $+0.27$ and $-0.33$ levels. For the portion the curriculum covers, the
standard is internally consistent.

\paragraph{The writable ceiling is invariant.} It sits between 57.7\% and
62.5\% on every shelf; across the corpus, median 60.1\% (IQR 55.2--64.9\%). It
does not vary with difficulty. Whether a text is a 49-character beginner story
or a 300-character literary passage, \textbf{about two characters in five are
ones no HSK level asks a learner to write by hand}. That is a property of the
standard, not of the text.

The allocation shows why. The syllabus issues writing characters as 100 across
levels 1--2 combined, then 150 per level, then 500 across 7--9 --- there is no
separate level-2 writing allocation. Of the 125 characters HSK 2 newly requires
a learner to \emph{read}, five are also required for writing. The two curricula
advance on different schedules, and level 2 advances only one.

\section{The corpus}
\label{sec:corpus}

We release 102 word-aligned graded texts: 1{,}185 sentences, 8{,}682 tokens,
14{,}417 graded characters, on six shelves, mean length rising from 49 to about
290 characters. Every token carries hanzi, pinyin and gloss.

The texts were \textbf{drafted with large-language-model assistance and then
reviewed, edited and re-levelled by the author}; the datasheet discloses this
fully. They are pedagogical material written to a level target, not naturally
occurring Chinese.

\section{HSKBench}
\label{sec:bench}

\subsection{Why not level prediction}

The obvious benchmark is level prediction. When gold labels come from the
grader under test, this measures only whether a model memorised a character
list: circular and uninformative.

We invert the relationship. \textbf{The grader is the metric, not the label.}
The task is generation: produce a coherent passage on a topic, at a length,
readable at a target level. Scoring runs the grader over the output as a
compiler is run over generated code --- deterministic, no judge model, no API.

\subsection{Task and metrics}

150 tasks: six target levels $\times$ 25 topics, five topics held out. Primary
metric is \textbf{level accuracy}, the share of outputs whose measured level
does not exceed the target. We also report mean signed error, length
compliance, and single-character budget violations. Unanswered tasks score as
failures. Coherence is excluded from the automatic score: judging it needs a
human or a model and would make results unreproducible.

\subsection{A human reference point}

\begin{table}[t]
\centering\small
\begin{tabular}{@{}lrr@{}}
\toprule
\textbf{System} & \textbf{Level acc.} & \textbf{Signed err.} \\
\midrule
\texttt{deepseek-chat} & \textbf{66.7\%} & +0.24 \\
\texttt{deepseek-reasoner} & 64.0\% & +0.49 \\
Human authors & 61.8\% & +0.27 \\
Corpus retrieval & 59.3\% & +0.27 \\
\bottomrule
\end{tabular}
\caption{Overall results. We show below that this table is close to
uninformative.}
\label{tab:bench}
\end{table}

\begin{table}[t]
\centering\small
\begin{tabular}{@{}lrrrrrr@{}}
\toprule
\textbf{Target} & 1 & 2 & 3 & 4 & 5 & 6 \\
\midrule
\texttt{chat}     & 16 & 24 & 72 & 88 & \textbf{100} & \textbf{100} \\
\texttt{reasoner} & \textbf{48} & \textbf{48} & \textbf{80} & 80 & 76 & 52 \\
Difference & +32 & +24 & +8 & $-$8 & $-$24 & $-$48 \\
\bottomrule
\end{tabular}
\caption{Level accuracy (\%) by target, both models at temperature 0.}
\label{tab:perlevel}
\end{table}

\subsection{The aggregate score is misleading}

Perfect scores at HSK 5--6 are not evidence of skill. At those levels the
constraint is nearly vacuous: almost any fluent Chinese passage satisfies a
95\% coverage bar drawn against the 1,940 characters graded at HSK 6 or below.
The constraint only
begins to bite at HSK 3, and by HSK 1 \texttt{deepseek-chat} fails 21 of 25
times. A system answering only the top three levels would score 96\%, so
\textbf{the aggregate must not be reported alone}.

\subsection{Reasoning helps at a floor and hurts at a ceiling}

\texttt{deepseek-reasoner} scores 2.7 points worse overall, and on the
aggregate is a slightly inferior system. Per level the two are barely the same
kind of model: the difference row in Table~\ref{tab:perlevel} is monotonic
across the whole scale, from $+32$ to $-48$. Averaged over the bottom half,
reasoning is worth \textbf{$+21.3$ points}; over the top half it costs
\textbf{$-26.7$}.

These are two different tasks wearing one instruction. Writing at HSK 1 is
constraint satisfaction --- stay inside a 246-character inventory --- which
rewards planning, and the reasoner's length compliance is also higher (80.0\%
against 73.3\%). Writing at HSK 6 is not a constraint problem at all, since
almost any fluent Chinese satisfies it; there deliberation hurts. Asked for
advanced prose, the reasoner reaches for advanced vocabulary and overshoots
into the 7--9 band, with mean signed error $+0.20$ at HSK 6 where
\texttt{deepseek-chat} sits at $-0.56$.

\textbf{Reasoning buys a floor and costs a ceiling.} Any evaluation of
difficulty-controlled generation that reports one number is measuring the wrong
thing.

Human accuracy fails \emph{directionally}: authors overshoot by 1.23 levels on
the easiest shelf and undershoot by 0.75 on the hardest, regressing toward
middle difficulty regardless of instruction. Writing to a level target is a
genuinely hard control problem for people, which is what makes the benchmark
non-trivial.

\section{Related work}

\paragraph{The documents' own literature.} \citet{feng2026hsk30} give the
authoritative account of the syllabus's development. It is qualitative: it
reports a high overlap rate and states that the syllabus was benchmarked
against the 2021 standard, but does not quantify how many words changed level,
nor compare the documents word by word. We are not aware of a published
quantitative comparison of the kind in \S\ref{sec:disagree}, in either
language.

We searched in both English and Chinese and found only adjacent prior work:
textbook-to-standard vocabulary alignment studies, and comparisons of the 1992
and 2009 HSK outlines. We were unable to search CNKI full text, so we state the
novelty claim as bounded rather than absolute: no quantitative comparison of
these two documents appears in the sources available to us, and given the
syllabus is under a year old, we consider a prior one unlikely but cannot
exclude it.

\paragraph{The problem is live in deployed tools.} Commercial HSK text checkers
grade against ``the new HSK 3.0 standard'' without naming which document, at
least one using greedy longest-match word segmentation --- the approach
\S\ref{sec:charlevel} shows to be unusable --- while publishing no data
provenance or version.

\paragraph{Chinese readability assessment.} The dominant approach is supervised
classification over engineered features. CRIE \citep{sung2016crie} extracts 82
multilevel features, trained on Taiwanese school textbooks; recent work applies
pretrained transformers with feature fusion \citep{yang2025chinese}. Our method
is deliberately not in this tradition: it is a transparent, deterministic
function of an official published standard, with no training data and no
learned parameters, and is auditable in a way a classifier is not.

\paragraph{Difficulty-controlled generation.} Controlling generated readability
is active, largely in English and against CEFR.
\citet{imperial2023flesch} evaluate instruction-tuned models against FKGL and
CEFR on teaching tasks. \citet{barayan2025analyzing} analyse zero-shot
readability-controlled simplification and report a limited ability to control
readability reliably. Our \S\ref{sec:bench} result is the same phenomenon in
Chinese, and sharpens it: control does not fail uniformly but collapses at one
end of the scale, precisely where the constraint binds. Our contribution is
also narrower and more reproducible, because the target is an explicit
character inventory rather than a learned readability model, so compliance is
checkable exactly and without a judge.

\section{Limitations}

Coverage is not comprehension: 95\% character coverage is necessary, not
sufficient, and grammar, register and world knowledge are unmodelled. The
threshold itself is imported from word-level English research; establishing it
empirically for L2 Chinese characters would be a contribution in itself. The
corpus is 102 texts --- adequate as a reference and validation set, inadequate
for training --- and the regrading result rests on it. It is LLM-assisted, as
disclosed. The 2021 lists come from public transcriptions of a scanned PDF and
the 2025 lists from our own extraction, validated against published totals but
not against an official machine-readable release, which does not exist. We
handle simplified characters only. Proper-noun exclusion needs pinyin, and
place names are a grey area: excluding 欧洲 is defensible under the
glossed-in-place principle and debatable under a vocabulary-knowledge one.
Shelf assignments are a single annotator's, with no inter-annotator agreement.

\section{Conclusion}

``HSK 3.0'' names two official documents that disagree on 41.5\% of their
shared vocabulary and 40.7\% of their shared characters, and choosing between
them changes the assigned level of nearly half of real texts, almost always
upward. The distinction is invisible in most current tooling. We release a
validated extraction of the examination syllabus, a tested grader that works
against either document and records which, an aligned corpus, and a benchmark
measuring the harder half of the problem.

\section*{Availability}

This paper is archived at \href{https://doi.org/10.5281/zenodo.22239032}{doi:10.5281/zenodo.22239032}.
Software, corpus and benchmark are archived separately at
\href{https://doi.org/10.5281/zenodo.22234657}{doi:10.5281/zenodo.22234657}
(concept DOI, resolving to the latest version). The library is on PyPI as
\texttt{hsk30}, the corpus on HuggingFace as
\texttt{harukicoder/hsk30-graded-readers}, and the source at
\url{https://github.com/harukicoder/hsk30}. Code is MIT; the corpus is CC BY
4.0. All figures regenerate with \texttt{scripts/reproduce.py}.

\section*{Ethics and data statement}

The corpus contains no personal data and describes invented characters in
everyday situations; its LLM-assisted authorship is disclosed in the datasheet.
Word and character lists derive from published national standards. No human
subjects were involved. An intended application is a Chinese-learning website
operated by the author.

% acl.sty already issues \bibliographystyle{acl_natbib}; declaring it again
% makes BibTeX abort with "Illegal, another \bibstyle command".
\bibliography{refs}

\end{document}
